\documentclass[10pt,a4paper]{amsart}
\usepackage[margin=19mm]{geometry}
\usepackage{amsmath,amssymb,mathtools}
\usepackage[scr=rsfs,cal=euler]{mathalfa}
\usepackage{xfrac}
\usepackage[unicode,hidelinks,bookmarks=false]{hyperref}
\newcommand{\ie}{i.e.\ }
\newcommand{\cf}{cf.\ }
\newcommand{\resp}{respectively\ }
\newcommand{\Subsection}{\subsection}
\providecommand{\nobreakdash}{\nobreak\hbox{-}}
\setlength{\emergencystretch}{2em}
\multlinegap0pt
\usepackage{amssymb}
\usepackage{dsfont}
\usepackage{xcolor}
\usepackage{enumerate}
\newtheorem{lemma}{Lemma}
\def\R{{\mathbb R}}% real numbers
\renewcommand{\geq}{\geqslant}\renewcommand{\ge}{\geqslant}
\renewcommand{\leq}{\leqslant}\renewcommand{\le}{\leqslant}
\title{Long title}
\author{Cassandre Lebot}
\date{Source: arXiv:2602.23189v3 (CC BY 4.0)}
\begin{document}
\maketitle

\noindent\emph{Source and licence.} Long title. Version arXiv:2602.23189v3 (CC BY 4.0), distributed by arXiv under the Creative Commons Attribution 4.0 International licence, http://creativecommons.org/licenses/by/4.0/, as recorded in the arXiv licence metadata for this exact version and shown on the abstract page https://arxiv.org/abs/2602.23189v3. Full licence notice: \texttt{licenses/arxiv-2602.23189-CC-BY-4.0.txt}.\par\smallskip

\noindent\textit{Editorial scope.} Counted unit: the article's lemma technicalemma (source0.tex lines 1119--1129). The article's complete original proof of it is delivered with it (source0.tex lines 1133--1190, one \texttt{proof} environment of 58 lines, which the article places away from the statement). No mathematics is written, repaired or re-derived: the statement and the proof are the article's own lines, in the article's own notation, and no retained line is substituted or re-flowed.  No further source-local context is needed: every cross-reference of every retained line resolves inside the two delivered spans.  Nothing was omitted from any retained span, so the package carries no omission record.  No retained line contains a citation, so the wrapper carries no bibliography and no bibliography entry.  No retained line contains a figure environment, an image-inclusion command or drawing code, so no image is delivered and the package declares no asset dependency.  Editorial formatting only, in addition: an AMS \texttt{amsart} preamble that restates the article's own declarations and macros used by the retained lines, namely the environments \texttt{lemma} on separate counters, together with the standard packages those lines need.  The retained environments print in this document as Lemma 1 in order of appearance, whereas the article prints the same items with the numbering of its own full document.  The complete native source of the article as distributed in the arXiv e-print for this exact version is bundled unchanged as \texttt{sources/195381/source0.tex}.
\begin{lemma}\label{technicalemma}
Let $\gamma >0$. For all $d_1, d_2 >0$ such that $d_2 \in [\underline{c},\overline{c}]\subset \R^+_*$, there exists a unique solution $(a_1,a_2) \in (0,1)$ that satisfies
\begin{align*}
 d_1 a_1^\gamma + (a_1-1) = 0,\\
 d_2 a_2^\gamma + (a_2-1) = 0.
\end{align*}
Moreover, there exists a constant $C$ (that only depends on $\gamma, \overline{c}, \underline{c}$) such that
\begin{equation*}
 \lvert a_1 - a_2 \rvert \leq C \lvert d_1 - d_2 \rvert.
\end{equation*}
\end{lemma}

\begin{proof}[Proof of Lemma~\ref{technicalemma}]
For each $i =1,2$, we define the function $F_{d_i} : \mathbb{R}_+ \to \mathbb{R}$ by
\[
F_{d_i}(x) := d_i x^\gamma + x - 1.
\]
The function $F_{d_i}$ is strictly increasing with respect to $x$. Moreover, we have
$F_{d_i}(0) = -1$ and $F_{d_i}(1) = d_i > 0$.
By continuity, there exists a unique solution $x_i \in (0,1)$ to the equation
$F_{d_i}(x) = 0$.
We denote this solution by $a_i$ for $i = 1,2$.

We then define the function $g:[\underline{c},\overline{c}]\to (0,1)$ by setting $a = g(d)$ such that $F_d(a)=0$. The monotonicity of $F_{d}$ with respect to $x$ implies that $g$ is decreasing with respect to $d$.

First, consider the case where $d_1 < \underline{c}/2$ or $d_1 > 2\overline{c}$. Since $d_2 \in [\underline{c}, \overline{c}]$, there exists a constant $C > 0$ sufficiently large such that
\[
|d_1 - d_2| \geq 1/C.
\]
Moreover, since $a_1, a_2 \in (0,1)$, it follows that
\[
|a_1 - a_2| \leq 1 \leq C |d_1 - d_2|.
\]

Now consider the case where $d_1 \in [\underline{c}/2,2 \overline{c}]$. Since the function $g$ is decreasing, it follows that $a_1 \in [g(2 \overline{c}), g(\frac{\underline{c}}{2})]$. We then introduce the compact set $K = [\frac{\underline{c}}{2},2 \overline{c}] \times [g(2 \overline{c}), g(\frac{\underline{c}}{2})]$ and and define the function $F: K \to \mathbb{R}$ by $F(d,x) = d x^\gamma +x-1$. The function $F$ is of class $C^\infty$ on $K$.

For any $(d_0,x_0) \in K$ satisfying $F(d_0,x_0)=0$, as $\partial_x F(d_0,x_0) \neq 0$, the Implicit Function Theorem ensures the existence of an open neighbourhood $U_{0,\mathrm{loc}}$ of $d_0$ and a unique function $g_{0,\mathrm{loc}} : U_{0,\mathrm{loc}} \to \mathbb{R}$ such that
\[
x = g_{0,\mathrm{loc}}(d)
\quad \text{and} \quad
F\bigl(d, g_{0,\mathrm{loc}}(d)\bigr)=0
\quad \text{for all } d \in U_{0,\mathrm{loc}}.
\]
Since the solution $(d,x)$ of the equation $F(d,x)=0$ is unique, it follows that
\[
g_{0,\mathrm{loc}} = g.
\]


From the Implicit Function Theorem, we also have
\[
g_{0,\mathrm{loc}}'(d) = -\,\frac{\partial_d F(d,g(d))}{\partial_x F(d,g(d))} = -\,\frac{x^\gamma}{\gamma d x^{\gamma-1}+1}, \qquad \text{with } x = g(d).
\]
Since $(d,x) \in K$, there exists a constant $L>0$ such that
\[
\left| g'(d) \right| \le L \quad \text{for all } d \in \Bigl[\frac{\underline{c}}{2},2 \overline{c}\Bigr].
\]
It then follows that
\begin{equation*}
 |a_1 - a_2| = |g(d_1) - g(d_2)| \le L\, |d_1 - d_2|.
\end{equation*}

Hence, we have shown that there exists a constant $C>0$ such that
\[
|a_1 - a_2| \le C\, |d_1 - d_2|,
\]
which concludes the proof.


\end{proof}

\end{document}
